\chapter{Szeregi liczbowe}
\label{ma1:ch:numerical-series}

Skończoną sumę liczb można obliczyć za pomocą działań w ciele. Suma nieskończenie wielu liczb wymaga nowej definicji: nie jest pojedynczym działaniem algebraicznym, lecz granicą ciągu sum częściowych. Rozwiniemy teorię jednocześnie dla liczb rzeczywistych i zespolonych; kryteria porównawcze wykorzystujące porządek sformułujemy dla ciągów nieujemnych. Zupełność przestrzeni $\R$ i $\C$ będzie istotna przy kryterium Cauchy'ego.

\section{Definicja szeregu i jego sumy}
\begin{definicja}
\label{ma1:def:series}
Niech $K=\R$ lub $K=\C$, a $(a_n)_{n\ge1}$ będzie ciągiem w $K$. Szeregiem o wyrazach $a_n$ nazywamy wyrażenie $\sum_{n=1}^\infty a_n$ rozpatrywane wraz z ciągiem sum częściowych
\[
 S_0=0,\qquad S_N=\sum_{n=1}^N a_n\quad(N\in\N_+).
\]
Szereg jest zbieżny, jeżeli istnieje $S\in K$ takie, że $S_N\to S$. Wówczas $S$ nazywamy sumą szeregu i piszemy $\sum_{n=1}^\infty a_n=S$. Jeżeli ciąg sum częściowych nie ma granicy w $K$, szereg jest rozbieżny.
\end{definicja}
Należy odróżniać wyraz szeregu $a_n$ od sumy częściowej $S_n$. Zbieżność ciągu $(a_n)$ nie oznacza zbieżności szeregu.

\begin{twierdzenie}[Warunek konieczny zbieżności]
\label{ma1:thm:series-term-test}
Jeżeli szereg $\sum_{n=1}^\infty a_n$ w $\R$ lub $\C$ jest zbieżny, to $a_n\to0$.
\end{twierdzenie}
\begin{proof}
Jeśli $S_n\to S$, to również $S_{n-1}\to S$: dla danego $\varepsilon>0$ nierówność $|S_j-S|<\varepsilon$ dla wszystkich $j\ge N$ daje ją dla $j=n-1$ przy $n\ge N+1$. Ponieważ $a_n=S_n-S_{n-1}$, rachunek granic (twierdzenie \ref{ma1:thm:limit-arithmetic}) daje $a_n\to S-S=0$.
\end{proof}

\begin{przyklad}
\label{ma1:ex:harmonic-series}
Szereg harmoniczny $\sum_{n=1}^\infty 1/n$ jest rozbieżny, mimo że $1/n\to0$. Istotnie, dla każdego $N\ge1$ jego sumy częściowe spełniają
\[
 S_{2N}-S_N=\sum_{n=N+1}^{2N}\frac1n\ge\frac12.
\]
Gdyby $(S_N)$ był zbieżny, byłby ciągiem Cauchy'ego na mocy twierdzenia \ref{ma1:thm:cauchy-basic}; powyższa nierówność przeczy temu dla $\varepsilon=1/2$. Warunek konieczny nie jest wystarczający.
\end{przyklad}

\begin{twierdzenie}[Kryterium Cauchy'ego dla szeregów]
\label{ma1:thm:series-cauchy}
Niech $(a_n)$ będzie ciągiem w $K=\R$ lub $K=\C$. Szereg $\sum_{n=1}^\infty a_n$ jest zbieżny wtedy i tylko wtedy, gdy
\[
 \forall\varepsilon\in\R_{>0}\quad\exists N\in\N_+\quad
 \forall p,q\in\N_+\quad
 \left(q\ge p\ge N\Rightarrow
 \left|\sum_{n=p}^{q}a_n\right|<\varepsilon\right).
\]
\end{twierdzenie}
\begin{proof}
Dla $q\ge p\ge1$ zachodzi $\sum_{n=p}^q a_n=S_q-S_{p-1}$. Jeśli szereg jest zbieżny, ciąg $(S_n)$ spełnia warunek Cauchy'ego. Dla $\varepsilon>0$ wybieramy indeks $M$ z $|S_i-S_j|<\varepsilon$ dla wszystkich $i,j\ge M$. Wystarczy w tezie przyjąć $N=M+1$, bo wtedy $q,p-1\ge M$.
Odwrotnie, załóżmy warunek z tezy. Dla $\varepsilon>0$ wybieramy odpowiednie $N$. Gdy $i>j\ge N$, mamy $|S_i-S_j|=|\sum_{n=j+1}^i a_n|<\varepsilon$, gdyż $i\ge j+1\ge N$. Gdy $i<j$ korzystamy z symetrii modułu, a dla $i=j$ różnica jest zerowa. Ciąg sum częściowych jest więc ciągiem Cauchy'ego. Zupełność $\R$ lub $\C$ (twierdzenie \ref{ma1:thm:real-metric-complete} i wniosek \ref{ma1:cor:euclidean-complete}) daje jego zbieżność do elementu $K$.
\end{proof}

\begin{twierdzenie}
\label{ma1:thm:series-linear-tail}
Niech $(a_n),(b_n)$ będą ciągami w $K=\R$ lub $K=\C$.
\begin{enumerate}[label=(\roman*)]
\item Dla każdego $m\in\N_+$ szereg $\sum_{n=1}^\infty a_n$ jest zbieżny wtedy i tylko wtedy, gdy zbieżny jest szereg $\sum_{n=m}^\infty a_n$. W przypadku zbieżności
$\sum_{n=1}^\infty a_n=\sum_{n=1}^{m-1}a_n+\sum_{n=m}^\infty a_n$.
\item Jeśli oba szeregi o wyrazach $a_n,b_n$ są zbieżne i $\lambda,\mu\in K$, to szereg o wyrazach $\lambda a_n+\mu b_n$ jest zbieżny oraz
\[
 \sum_{n=1}^\infty(\lambda a_n+\mu b_n)
 =\lambda\sum_{n=1}^\infty a_n+\mu\sum_{n=1}^\infty b_n.
\]
\end{enumerate}
\end{twierdzenie}
\begin{proof}
Sumę szeregu zaczynającego się od $m$ definiujemy jako granicę $\sum_{n=m}^N a_n$ dla $N\to\infty$, $N\ge m$. Dla (i) sumy te są równe $S_N-S_{m-1}$, gdzie $S_0=0$. Dodanie lub odjęcie stałej zachowuje zbieżność, na mocy twierdzenia \ref{ma1:thm:limit-arithmetic}; przejście do granicy daje wzór. Dla (ii) suma częściowa po lewej jest równa $\lambda S_N+\mu T_N$, gdzie $S_N,T_N$ są sumami częściowymi obu szeregów. Rachunek granic daje tezę.
\end{proof}

\section{Szeregi geometryczne i teleskopowe}
\begin{twierdzenie}[Szereg geometryczny]
\label{ma1:thm:geometric-series}
Niech $q\in\C$. Szereg $\sum_{n=0}^\infty q^n$, gdzie $q^0=1$ także dla $q=0$, jest zbieżny wtedy i tylko wtedy, gdy $|q|<1$. W przypadku zbieżności jego suma wynosi $1/(1-q)$.
\end{twierdzenie}
\begin{proof}
Dla $q\ne1$ i $N\ge0$ z rozdzielności mnożenia mamy
\[
 (1-q)\sum_{n=0}^N q^n=\sum_{n=0}^N q^n-\sum_{n=1}^{N+1}q^n
 =1-q^{N+1}.
\]
Stąd suma częściowa jest równa $(1-q^{N+1})/(1-q)$. Dla $|q|<1$ z przykładu \ref{ma1:ex:geometric-sequence} wynika $q^{N+1}\to0$, więc suma częściowa dąży do $1/(1-q)$. Jeżeli $|q|\ge1$, mamy $|q^n|\ge1$ dla wszystkich $n\ge1$, co przeczy warunkowi koniecznemu zbieżności z twierdzenia \ref{ma1:thm:series-term-test}. Obejmuje to także $q=1$.
\end{proof}

\begin{przyklad}
Dla $q=1/2$ suma częściowa od $n=0$ do $n=N$ wynosi $2-2^{-N}$, a suma szeregu $2$. Błąd po obliczeniu tej sumy częściowej wynosi dokładnie $2^{-N}$. Dla dowolnego $\varepsilon>0$ można dobrać $N$ tak, aby $2^{-N}<\varepsilon$, bez obliczania nieskończonej sumy.
\end{przyklad}

\begin{twierdzenie}[Szereg teleskopowy]
\label{ma1:thm:telescoping-series}
Niech $(b_n)_{n\ge1}$ będzie ciągiem w $K=\R$ lub $K=\C$. Szereg $\sum_{n=1}^\infty(b_n-b_{n+1})$ jest zbieżny wtedy i tylko wtedy, gdy ciąg $(b_n)$ jest zbieżny. Jeżeli $b_n\to b$, to suma szeregu wynosi $b_1-b$.
\end{twierdzenie}
\begin{proof}
Dla $N\ge1$ mamy
\[
 \sum_{n=1}^N(b_n-b_{n+1})=b_1-b_{N+1},
\]
ponieważ każdy wyraz $b_j$ dla $2\le j\le N$ występuje dwukrotnie, z przeciwnymi znakami. Zbieżność $(b_{N+1})$ jest równoważna zbieżności $(b_n)$ do tej samej granicy, gdyż pominięcie pierwszego wyrazu nie wpływa na definicję granicy. Rachunek granic daje obie implikacje i wzór.
\end{proof}

\begin{przyklad}
\label{ma1:ex:telescoping-series}
Z tożsamości $1/(n(n+1))=1/n-1/(n+1)$ otrzymujemy
\[
 \sum_{n=1}^N\frac1{n(n+1)}=1-\frac1{N+1},
 \qquad \sum_{n=1}^\infty\frac1{n(n+1)}=1.
\]
Różnica sumy szeregu i sumy częściowej wynosi $1/(N+1)$.
\end{przyklad}

\section{Szeregi o wyrazach nieujemnych}
\begin{twierdzenie}
\label{ma1:thm:positive-series}
Niech $a_n\ge0$ dla każdego $n\ge1$, gdzie $a_n\in\R$. Szereg $\sum_{n=1}^\infty a_n$ jest zbieżny wtedy i tylko wtedy, gdy jego sumy częściowe są ograniczone z góry. W przypadku zbieżności
\[
 \sum_{n=1}^\infty a_n=\sup_{N\ge1}\sum_{n=1}^N a_n.
\]
Jeśli sumy częściowe nie są ograniczone z góry, to $S_N\to+\infty$.
\end{twierdzenie}
\begin{proof}
Ponieważ $S_{N+1}-S_N=a_{N+1}\ge0$, ciąg $(S_N)$ jest niemalejący. Jeśli jest ograniczony z góry, twierdzenie \ref{ma1:thm:monotone-convergence} daje zbieżność do kresu górnego. Odwrotnie, zbieżny ciąg jest ograniczony na mocy twierdzenia \ref{ma1:thm:convergent-bounded}. Gdy $(S_N)$ jest nieograniczony z góry, dla każdego $M\in\R$ istnieje $N$ z $S_N>M$, a monotoniczność daje $S_n\ge S_N>M$ dla każdego $n\ge N$.
\end{proof}

\begin{twierdzenie}[Kryterium porównawcze]
\label{ma1:thm:series-comparison}
Niech $(a_n),(b_n)$ będą ciągami rzeczywistymi i niech od pewnego indeksu $0\le a_n\le b_n$. Jeżeli $\sum b_n$ jest zbieżny, to $\sum a_n$ jest zbieżny. Jeżeli $\sum a_n$ jest rozbieżny, to $\sum b_n$ jest rozbieżny.
\end{twierdzenie}
\begin{proof}
Wybieramy $m$ tak, aby nierówności zachodziły dla $n\ge m$. Z twierdzenia \ref{ma1:thm:series-linear-tail}(i) wystarczy badać szeregi od $m$. Dla $N\ge m$ zachodzi $0\le\sum_{n=m}^N a_n\le\sum_{n=m}^N b_n$. Jeśli drugi szereg jest zbieżny, jego sumy częściowe są ograniczone z góry. Pierwsze sumy częściowe również są ograniczone z góry i niemalejące, więc są zbieżne na mocy twierdzenia \ref{ma1:thm:positive-series} po przeindeksowaniu. Druga teza jest kontrapozycją pierwszej.
\end{proof}

\begin{wniosek}[Porównanie graniczne]
\label{ma1:cor:limit-comparison}
Niech $a_n\ge0$ i $b_n>0$ od pewnego indeksu, gdzie $(a_n),(b_n)$ są ciągami rzeczywistymi. Jeśli $a_n/b_n\to c\in(0,+\infty)$, to oba szeregi $\sum a_n$ i $\sum b_n$ są jednocześnie zbieżne albo jednocześnie rozbieżne. Jeśli $a_n/b_n\to0$, zbieżność $\sum b_n$ implikuje zbieżność $\sum a_n$. Jeśli $a_n/b_n\to+\infty$, rozbieżność $\sum b_n$ implikuje rozbieżność $\sum a_n$.
\end{wniosek}
\begin{proof}
W pierwszym przypadku od pewnego indeksu $c/2<a_n/b_n<3c/2$, więc $(c/2)b_n\le a_n\le(3c/2)b_n$. Zbieżność szeregu o wyrazach $\lambda b_n$ dla $\lambda>0$ jest równoważna zbieżności szeregu o wyrazach $b_n$, ponieważ można pomnożyć sumy częściowe przez $\lambda$ lub $1/\lambda$. Dwukrotne użycie kryterium porównawczego daje równoważność.
Dla $c=0$ od pewnego indeksu $0\le a_n/b_n<1$, czyli $a_n\le b_n$. Dla granicy $+\infty$ od pewnego indeksu $a_n/b_n>1$, czyli $b_n\le a_n$. Stosujemy odpowiedni kierunek kryterium porównawczego.
\end{proof}

\section{Kryterium zagęszczania}
\begin{twierdzenie}[Kryterium Cauchy'ego przez zagęszczanie]
\label{ma1:thm:condensation}
Niech $(a_n)$ będzie nierosnącym ciągiem nieujemnych liczb rzeczywistych. Szereg $\sum_{n=1}^\infty a_n$ jest zbieżny wtedy i tylko wtedy, gdy zbieżny jest szereg $\sum_{k=0}^\infty 2^k a_{2^k}$.
\end{twierdzenie}
\begin{proof}
Dla każdego całkowitego $k\ge0$ blok indeksów $2^k\le n\le2^{k+1}-1$ ma $2^k$ elementów. Z monotoniczności mamy
\[
 2^k a_{2^{k+1}}\le
 \sum_{n=2^k}^{2^{k+1}-1}a_n
 \le 2^k a_{2^k}.
\]
Sumując po $k=0,\ldots,m$ dla dowolnego $m\ge0$, otrzymujemy
\[
 \frac12\sum_{j=1}^{m+1}2^j a_{2^j}
 \le S_{2^{m+1}-1}
 \le\sum_{k=0}^{m}2^k a_{2^k}.
\]
Jeśli szereg wyjściowy jest zbieżny, wszystkie $S_N$ są ograniczone, a lewe oszacowanie ogranicza sumy częściowe szeregu zagęszczonego od $j=1$. Dodanie skończonego wyrazu dla $j=0$ zachowuje ograniczoność. Twierdzenie \ref{ma1:thm:positive-series} daje zbieżność szeregu zagęszczonego.
Jeśli szereg zagęszczony jest zbieżny, prawe oszacowanie ogranicza $S_{2^{m+1}-1}$ niezależnie od $m$. Dla każdego $N$ istnieje $m$ z $2^{m+1}-1\ge N$ (np. $m=N$, ponieważ $2^{N+1}\ge N+2$ przez indukcję). Nieujemność wyrazów daje $S_N\le S_{2^{m+1}-1}$, więc wszystkie sumy częściowe są ograniczone z góry. Ponownie stosujemy twierdzenie \ref{ma1:thm:positive-series}.
\end{proof}

\begin{przyklad}
\label{ma1:ex:integer-p-series}
Dla dodatniej liczby całkowitej $p$ ciąg $a_n=1/n^p$ jest nierosnący i dodatni. Szereg zagęszczony ma wyrazy $2^k/(2^k)^p=(2^{1-p})^k$. Z twierdzenia \ref{ma1:thm:geometric-series} jest on zbieżny dokładnie wtedy, gdy $2^{1-p}<1$, czyli $p>1$. Stąd $\sum 1/n^p$ jest zbieżny dla całkowitych $p\ge2$, a dla $p=1$ rozbieżny. Przypadek dowolnego rzeczywistego wykładnika $p$ omówimy po skonstruowaniu potęg o wykładnikach rzeczywistych.
\end{przyklad}

\section{Zbieżność bezwzględna}
\begin{definicja}
\label{ma1:def:absolute-series}
Niech $(a_n)$ będzie ciągiem w $K=\R$ lub $K=\C$. Szereg $\sum a_n$ nazywamy zbieżnym bezwzględnie, jeśli szereg nieujemnych liczb rzeczywistych $\sum |a_n|$ jest zbieżny.
\end{definicja}
\begin{twierdzenie}
\label{ma1:thm:absolute-series}
Każdy szereg bezwzględnie zbieżny w $\R$ lub $\C$ jest zbieżny. Ponadto dla każdego $N\ge0$ zachodzi
\[
 \left|\sum_{n=1}^\infty a_n-S_N\right|\le
 \sum_{n=N+1}^\infty|a_n|.
\]
\end{twierdzenie}
\begin{proof}
Dla $\varepsilon>0$ z kryterium Cauchy'ego zastosowanego do $\sum|a_n|$ istnieje $m\in\N_+$ takie, że dla $q\ge p\ge m$ zachodzi $\sum_{n=p}^q|a_n|<\varepsilon$. Z nierówności trójkąta stosowanej indukcyjnie mamy
$|\sum_{n=p}^q a_n|\le\sum_{n=p}^q|a_n|<\varepsilon$.
Kryterium \ref{ma1:thm:series-cauchy} daje zbieżność szeregu $\sum a_n$. Niech $S$ będzie jego sumą. Dla $M>N$ zachodzi
$|S_M-S_N|\le\sum_{n=N+1}^M|a_n|$. Przy ustalonym $N$ prawa strona zbiega do $\sum_{n=N+1}^\infty|a_n|$, a lewa do $|S-S_N|$ na mocy twierdzenia \ref{ma1:thm:distance-limit}. Zachowanie nierówności nieostrych daje oszacowanie.
\end{proof}
\begin{przyklad}
Jeżeli $|a_n|\le Cq^n$ dla wszystkich $n\ge1$, gdzie $C\ge0$ oraz $0\le q<1$ są rzeczywiste, to szereg jest zbieżny bezwzględnie przez porównanie z szeregiem geometrycznym. Dla $N\ge0$ jego reszta spełnia
\[
 |S-S_N|\le C\sum_{n=N+1}^\infty q^n
 =\frac{Cq^{N+1}}{1-q}.
\]
Wzór obowiązuje również dla $q=0$, kiedy wszystkie $a_n=0$. Oszacowanie daje konkretną kontrolę błędu obcięcia szeregu.
\end{przyklad}

\begin{stwierdzenie}
\label{ma1:prop:complex-series-components}
Niech $a_n=x_n+\mathrm i y_n$, gdzie $x_n,y_n\in\R$. Szereg $\sum a_n$ jest zbieżny wtedy i tylko wtedy, gdy oba szeregi $\sum x_n$ i $\sum y_n$ są zbieżne, a wtedy
\[
 \sum_{n=1}^\infty a_n=\sum_{n=1}^\infty x_n+
 \mathrm i\sum_{n=1}^\infty y_n.
\]
Zbieżność bezwzględna $\sum a_n$ jest równoważna jednoczesnej zbieżności $\sum|x_n|$ i $\sum|y_n|$.
\end{stwierdzenie}
\begin{proof}
Części rzeczywista i urojona sumy częściowej są odpowiednimi sumami częściowymi szeregów rzeczywistych. Kryterium zbieżności współrzędnych w $\C$ z rozdziału \ref{ma1:ch:metric-sequences} daje pierwszą równoważność i wzór na sumę. Dla drugiej korzystamy z $|x_n|,|y_n|\le|a_n|\le|x_n|+|y_n|$. Jeśli $\sum|a_n|$ jest zbieżny, kryterium porównawcze daje zbieżność obu szeregów współrzędnych. Jeśli oba są zbieżne, to ich suma jest zbieżna z twierdzenia \ref{ma1:thm:series-linear-tail}, a kryterium porównawcze daje zbieżność $\sum|a_n|$.
\end{proof}

\section{Kryteria ilorazowe i pierwiastkowe}
\begin{twierdzenie}[Kryterium d'Alemberta]
\label{ma1:thm:ratio-test}
Niech $(a_n)$ będzie ciągiem w $\R$ lub $\C$ i niech $a_n\ne0$ od pewnego indeksu. Oznaczmy $r_n=|a_{n+1}|/|a_n|$ dla tych indeksów.
\begin{enumerate}[label=(\roman*)]
\item Jeśli $\limsup r_n<1$, szereg $\sum a_n$ jest zbieżny bezwzględnie.
\item Jeśli $\liminf r_n>1$, szereg $\sum a_n$ jest rozbieżny.
\end{enumerate}
Jeżeli istnieje granica $r=\lim r_n$ w rozszerzonej prostej rzeczywistej, wystarczy w (i) założyć $r<1$, a w (ii) $r>1$.
\end{twierdzenie}
\begin{proof}
Początkowe nieokreślone ilorazy możemy uzupełnić dowolnymi liczbami rzeczywistymi, gdyż nie wpływa to na granicę dolną ani górną. W przypadku (i) wybieramy rzeczywiste $q$ z $\limsup r_n<q<1$ oraz $q>0$. Z kryterium progowego (twierdzenie \ref{ma1:thm:limsup-thresholds}) istnieje $m$ takie, że dla każdego $n\ge m$ mamy $a_n\ne0$ i $r_n<q$. Indukcja daje dla $k\ge0$
$|a_{m+k}|\le |a_m|q^k$: dla $k=0$ zachodzi równość, a przejście do $k+1$ wynika z $|a_{m+k+1}|<q|a_{m+k}|$. Szereg modułów od $m$ jest zbieżny przez porównanie z szeregiem geometrycznym, a skończony początek nie wpływa na zbieżność.
W przypadku (ii) wybieramy rzeczywiste $q>1$ z $q<\liminf r_n$, również gdy granica dolna wynosi $+\infty$. Od pewnego indeksu $m$ mamy $r_n>q$, więc indukcyjnie $|a_{m+k}|\ge|a_m|q^k\ge|a_m|>0$. Warunek konieczny $a_n\to0$ nie zachodzi. Ostatnie zdanie wynika z twierdzenia \ref{ma1:thm:liminf-limsup-convergence}.
\end{proof}

\begin{twierdzenie}[Kryterium pierwiastkowe Cauchy'ego]
\label{ma1:thm:root-test}
Niech $(a_n)$ będzie ciągiem w $\R$ lub $\C$ i niech
$L=\limsup_{n\to\infty}\sqrt[n]{|a_n|}$, gdzie kres górny granic bierzemy w rozszerzonej prostej rzeczywistej.
\begin{enumerate}[label=(\roman*)]
\item Jeżeli $L<1$, szereg $\sum a_n$ jest zbieżny bezwzględnie.
\item Jeżeli $L>1$, szereg $\sum a_n$ jest rozbieżny.
\end{enumerate}
\end{twierdzenie}
\begin{proof}
Dla $L<1$ wybieramy rzeczywiste $q$ z $L<q<1$ i $q>0$. Z kryterium progowego od pewnego indeksu $\sqrt[n]{|a_n|}<q$. Ponieważ potęgowanie stopnia $n$ jest ściśle rosnące na liczbach nieujemnych, mamy $|a_n|<q^n$. Kryterium porównawcze daje zbieżność szeregu modułów.
Jeśli $L>1$, wybieramy $1<q<L$ (dla $L=+\infty$ dowolne rzeczywiste $q>1$). Dla każdego $N$ istnieje $n\ge N$ z $\sqrt[n]{|a_n|}>q$, więc $|a_n|>q^n\ge1$. To zaprzecza zbieżności $a_n$ do zera, zatem szereg jest rozbieżny.
\end{proof}
\begin{przyklad}
\label{ma1:ex:ratio-root-boundary}
Oba kryteria nie rozstrzygają przypadku granicy równej $1$. Dla $a_n=1/n$ i $b_n=1/n^2$ ilorazy kolejnych wyrazów są odpowiednio $n/(n+1)$ i $(n/(n+1))^2$, a więc oba zbiegają do $1$. Pierwszy szereg jest rozbieżny, drugi zbieżny.
Wykażmy również, że $\sqrt[n]{1/n}\to1$ oraz $\sqrt[n]{1/n^2}\to1$ bez używania funkcji wykładniczej. Dla dowolnego rzeczywistego $\delta>0$ i $n\ge3$ z rozwinięcia iloczynu $(1+\delta)^n$ otrzymujemy
\[
 (1+\delta)^n\ge \binom n3\delta^3
 =\frac{n(n-1)(n-2)}6\delta^3.
\]
Dla $n\ge4$ mamy $n-1\ge n/2$, $n-2\ge n/2$, stąd prawa strona jest nie mniejsza niż $n^3\delta^3/24$. Jeśli ponadto $n>24/\delta^3$, to $(1+\delta)^n>n^2$. Zatem od pewnego indeksu
$1\le\sqrt[n]{n}\le\sqrt[n]{n^2}<1+\delta$.
Oba ciągi pierwiastków dążą do $1$, a z rachunku granic ich odwrotności także. Granica górna w kryterium pierwiastkowym wynosi $1$ w obu przykładach.
\end{przyklad}
\begin{przyklad}
Kryterium pierwiastkowe nie wymaga niezerowych wyrazów. Niech $a_{2k}=2^{-2k}$ i $a_{2k-1}=0$ dla $k\ge1$. Wówczas $\sqrt[n]{|a_n|}$ przyjmuje wartości $1/2$ i $0$, obie nieskończenie wiele razy. Jego granica górna to $1/2<1$, więc szereg jest zbieżny bezwzględnie. Kryterium ilorazowe w podanej wersji nie ma zastosowania.
\end{przyklad}

\section{Przykłady kontroli błędu}
\begin{przyklad}
\label{ma1:ex:factorial-series}
Rozważmy szereg $\sum_{n=1}^\infty 1/n!$, gdzie $n!=1\cdot2\cdots n$. Iloraz kolejnych wyrazów wynosi $1/(n+1)$ i dąży do zera, więc szereg jest zbieżny. Dla $N\ge1$ i $k\ge1$ mamy
\[
 \frac1{(N+k)!}=\frac1{N!}\frac1{(N+1)\cdots(N+k)}
 \le\frac1{N!}\frac1{(N+1)^k}.
\]
Sumując dla $k\ge1$ i stosując kryterium porównawcze, otrzymujemy dla sumy $S$ szeregu
\[
 0\le S-\sum_{n=1}^N\frac1{n!}\le
 \frac1{N!}\frac{1/(N+1)}{1-1/(N+1)}
 =\frac1{NN!}.
\]
Przykład pokazuje, że oszacowania dowodzące zbieżności pozwalają też kontrolować dokładność obliczeń. Szereg ten będzie istotny przy konstrukcji funkcji wykładniczej; nie przyjmujemy jeszcze żadnej definicji tej funkcji ani jej własności.
\end{przyklad}
\begin{uwaga}
Zbieżność bezwzględna jest silniejsza od zwykłej zbieżności. W następnym rozdziale zbadamy szeregi naprzemienne, zbieżność warunkową oraz wpływ zmiany kolejności wyrazów. Nie zakładamy na tym etapie, że dowolne przestawienie nieskończenie wielu składników zachowuje sumę.
\end{uwaga}
